% Context from chapters/wavelets.tex; for mathematical reference, not compilation.
\section{On Bounded Domains}

On the periodic domain $\TT=\RR/\ZZ$, we use period 1 rather than the $2\pi$-period convention of the Fourier chapter. To obtain an orthonormal wavelet basis of $L^2(\TT)$, periodize the wavelets and restrict their translation points $2^j n$ to $[0,1]$, choosing $0 \leq n < 2^{-j}$. As in~\eqref{eq-periodizing}, the periodization of $f \in L^1(\RR)$ is
\eq{
	f^P = \sum_{n \in \ZZ} f(\cdot-n) \in L^1(\TT). 
}
For an integer coarsest level $j_0\leq0$, this gives the family
\eq{
	\enscond{\psi_{j,n}^P}{j \leq j_0, 0 \leq n < 2^{-j}} \cup 
	\enscond{\phi_{j_0,n}^P}{0 \leq n < 2^{-j_0}}.
}
which is an orthonormal basis of $L^2(\TT)$, illustrated in Figure~\ref{fig-periodize}.
 
One can also construct wavelet bases using Neumann (mirror) boundary conditions, but this is more involved.



\begin{figure}[tbp]
\centering
\BookPanel{.485}{1}{tikz/wavelets-periodization}\hfill
\BookPanel{.485}{2}{tikz/wavelets-periodic-translates}
\caption{Periodized basis functions. Only the translations $0\leq n<2^{-j}$ are distinct, and $j_0=0$ is a common choice of coarsest scale.}
\label{fig-periodize}
\end{figure}



                                                   
\section{Fast Wavelet Transform}\label{sec-fast-wavelet-transform}

   
\subsection{Discretization}

We now work over $\RR/\ZZ$.
 
We assume access to a discrete signal $a_J \in \RR^N$ with $N = 2^{-J}$ at a fixed scale $2^J$, whose entries are inner products with the scaling functions, i.e.
\eql{\label{eq-hyp-wav-samples}
	\foralls n\in \{0,\ldots,N-1\}, \quad a_{J,n} = \dotp{f}{\phi_{J,n}^P} \approx 2^{J/2}f(2^Jn),
}
for the underlying continuous function $f$. The approximation uses the normalization $\int\phi=1$. The equality assumes that acquisition gives exact access to $P_{V_J}f$, much as Shannon sampling assumes bandlimiting. This model is reasonable when the scaling functions $(\phi_{J,n})_n$ approximate the acquisition device's point-spread function. Preprocessing the measurements can sometimes improve agreement with~\eqref{eq-hyp-wav-samples}.

Starting from $a_J$, the discrete wavelet transform computes the coefficients
\eq{
	\foralls j \in \{J+1,J+2, \ldots, 0\}, \quad
	\foralls n \in \range{0,2^{-j}-1}, \quad
	a_{j,n} \eqdef \dotp{f}{\phi_{j,n}^P}, \qandq
	d_{j,n} \eqdef \dotp{f}{\psi_{j,n}^P}
}
in order of increasing $j$. After an approximation vector $a_j$ has been used to compute the next scale, it can be discarded; the detail vectors $d_j$ are retained.

The forward discrete wavelet transform on a bounded domain is thus the orthogonal finite-dimensional map
\eq{
	a_J \in \RR^N \longmapsto
	\enscond{ d_{j,n} }{ J<j\leq0,\ 0\leq n<2^{-j} } \cup \{a_0 \in \RR\}.
}
By orthogonality, the inverse transform is the adjoint.


Figure~\ref{fig-wavcoef-1d} shows examples of wavelet coefficients. For each scale $2^j$, there are $2^{-j}$ coefficients.

\myfigure{
\tabun{
\image{wavelets}{.9}{wavcoef-1d-all}\\
Wavelet coefficients $\{d_{j,n}\}_{j,n}$
}
\tabdeux{
\image{wavelets}{.48}{wavcoef-1d-0}&
\image{wavelets}{.48}{wavcoef-1d-1}\\
Signal $f$ & $d_{-7}$ \\
\image{wavelets}{.48}{wavcoef-1d-2}&
\image{wave